\documentclass{rapportCS}
\usepackage{lipsum}
\usepackage{pdfpages}
\usepackage{tikz}
\usepackage{microtype}
\usepackage[T1]{fontenc}
\usepackage{txfonts}
\usepackage{enumitem}

%----------- Crédits -------------------
% Template de rapport pour l'INSA Rennes basé sur le template "Template Stage Master IA - Université Le Mans" créé par Marie Tahon en 2023. 
% Thanks for Template Stage Master IA - Université Le Mans, By Marie Tahon (2023) that was the base for the creation of this template for INSA Rennes activity reports.

\title{Rapport Stage INSA Rennes 2023 - Template} 
\begin{document}

%----------- Informations du rapport ---------
\titreligneun{\uppercase{Titre du }}
\titrelignedeux{\uppercase{Document}}% Titre du document
\soustitre{Rapport de Projet de Fin d'Études} % Sous-titre du document pour mentionné le type de document
\numeroordre{Numéro d'ordre du PFE}
\departement{INFO} % Département pau sein de l'INSA
\logoentreprise{\includegraphics[scale=0.03]{logos/Google_2015_logo.png}} % Logo de l'entreprise pour inclure dans les marges du rapport, ajuster l'échelle selon la taille du logo.




%----------- Initialisation -------------------
\fairemarges % Afficher les marges
\fairecouverture % Créer la 1ère de couverture
% Autorisation de diffusion et d'archivage à ajouter en PDF.
\includepdf[pages=-,noautoscale = true,scale=1]{documents/ADA.pdf} % remplacez le fichier ADA.pdf par votre ADA remplie et signées par toutes les parties

%------------ Remerciements ----------------
\newpage
\pagenumbering{arabic} % Commencer la numérotation des pages
\thispagestyle{empty} % Exclure cette page des marges et des numéros
\section*{Remerciements}
    \lipsum[1] % Effacer la commande \lipsum et ajouter les remerciements à la place

%------------ Table des matières ----------------
\newpage
\tabledematieres % Créer la table de matières

%------------ Corps du rapport ----------------

% LES COMMANDES \lipsum[] SERVENT UNIQUEMENT DE REMPLISSAGE POUR ILLUSTRER ELLES DOIVENT ËTRE SUPPRIMEES

%------------ Introduction ----------------
\section{Introduction} 
% Effacer les lignes suivantes et écrire le texte souhaité
\lipsum[1]

\newpage
\section{Deuxième Section}
\subsection{Première Sous-section}
\subsubsection{Première Sous-sous-section}

\par Lorem ipsum dolor sit amet, consectetur adipiscing elit. Mauris semper bibendum finibus. Maecenas tristique feugiat luctus. Nunc vitae ante vitae magna venenatis mattis. Duis arcu nibh, ultrices vel auctor at, eleifend a libero. Donec in neque porttitor, venenatis nisl et, euismod dui. Donec sit amet rhoncus lectus, quis viverra metus. Nunc suscipit urna commodo, gravida ante a, pharetra elit. Etiam accumsan sed augue ac iaculis. Nunc et magna tellus. Duis in ante erat. Vivamus rhoncus, nisl a ultrices congue, enim magna elementum mauris, non laoreet nisl massa in velit. Sed congue mi ac velit elementum accumsan. Nulla facilisis blandit elementum. Nunc ex odio, vehicula ut lorem et, faucibus congue libero. Etiam vehicula elit a orci bibendum sodales. Integer viverra gravida magna, ut sollicitudin erat.

\par Pellentesque et pharetra libero. Praesent aliquam sapien vitae ligula lobortis, vitae dignissim ligula vestibulum. Donec aliquam dolor aliquet mi aliquet faucibus. Sed ullamcorper molestie dui eget euismod. Curabitur a orci sed ante posuere pharetra. Vivamus felis metus, cursus nec quam non, eleifend mattis risus. Fusce ornare luctus finibus. Nam maximus semper mollis. Mauris volutpat ultrices arcu, in bibendum tellus laoreet eu. Nullam condimentum ligula et risus consectetur, in tempus diam pharetra. Donec in eros at dolor porttitor venenatis sit amet et sapien. Morbi nec orci mauris. Donec porta, nibh in eleifend ullamcorper, ex dui efficitur sapien, vel aliquet elit dui at nisi. Ut sit amet fringilla diam, ut pharetra urna. Ut porttitor sagittis neque.

\par Donec a nisi ac lacus sodales scelerisque. Nam viverra tincidunt dapibus. Maecenas nisl felis, fringilla non ante vitae, tristique semper sem. Phasellus eu convallis lectus, non fermentum ipsum. Cras fringilla enim vitae nisi tempor, ut faucibus dolor sollicitudin. Aenean lacinia justo at ex faucibus gravida. Suspendisse sed eros eget diam consectetur ornare a in nisi. Quisque mauris urna, lacinia vel facilisis in, congue vel augue. Maecenas porttitor sit amet odio et hendrerit. Ut posuere vulputate nibh tempor tincidunt. Nunc ut est quis erat vehicula vestibulum vitae non lacus.

\subsubsection{Deuxième Sous-sous-section}
\lipsum[1]
\subsection{Deuxième Sous-section}
\lipsum[1]
\subsection{Troisième Sous-section}
\lipsum[1]

\section{Troisième Section}
\lipsum[1-3]

%------------- Commandes utiles ----------------
\newpage
\section{Quelques commandes}
\subsection{Insertion de listes}
\textbf{Consignes du moodle INFO, à adapter selon le département :}
\newline
Le rapport comprend nécessairement des parties incontournables :
\begin{itemize}[label=$\bullet$] %
    \item Une \textbf{couverture} respectant le modèle unifié imposé (disponible au format .doc) et mentionnant votre numéro d'ordre (visible sur la liste des correspondants enseignants). 
    \item La mise en page du contenu de votre document reste libre.
    \item Si vous rédigez en LaTeX vous pouvez rajouter la couverture (après conversion en .pdf du document obtenu à partir du fichier .doc) :
         \begin{enumerate}
            \item page 1 : Couverture avec le titre
            \item page 2-3 : ADA
            \item contenu édité en LaTeX
            \item page N-1 : double résumé 
            \item page N : dernière couverture (avec le logo INSA)
        \end{enumerate}

    \item Une partie introductive décrivant succinctement l'entreprise d'accueil et le sujet du stage (de quelques pages tout au plus).
    \item Une partie principale décrivant le travail réalisé pendant le stage. Il est utile de structurer ce travail, de hiérarchiser l'importance de chaque axe (si plusieurs) afin de mieux mettre en évidence les aspects techniques complexes, non triviaux (i.e., qui prouvent qu'il s'agit d'un travail de niveau M2), sans toutefois rentrer dans tous les détails techniques des aspects moins importants. La présentation doit être « didactique » : elle doit être compréhensible par un informaticien non spécialiste du domaine d'application de votre stage. Elle doit mettre en évidence le travail que vous avez effectivement réalisé.
    \item Une partie conclusion montrant d'une part, les apports de vos travaux pour l'entreprise et d'autre part, les enseignements que vous avez personnellement retirés de ce stage.
    \item Une partie annexe donnant un planning du stage en semaines indiquant la durée des différentes phases (étude, analyse fonctionnelle, codage, tests...).
    \item Une partie bibliographique comparant vos travaux avec ceux déjà effectués sera parfois nécessaire dans le document.
    \item Un double résumé français et anglais (une page au total, en 3e de couverture, c-à-d l'avant-dernière page du rapport).
        \begin{itemize}[label=$\diamond$]
            \item La consigne associée à ce résumé bilingue est d'être didactique et accessible à un non informaticien, pour par exemple être compris d'un jeune étudiant intéressé par le département informatique.
        \end{itemize}
\end{itemize}
Ce rapport doit être rédigé en français ou en anglais et ne pas dépasser 30 à 40 pages pour les parties mentionnées ci-dessus. Il peut être associé à une annexe technique comportant autant de documents que vous souhaitez absolument adjoindre au rapport final, mais dont la lecture doit demeurer facultative.
\subsection{Insertion de figures}
%------ Pour insérer et citer une image centralisée -----
% Le premier argument est le chemin pour la photo
% Le deuxième est la hauteur de la photo
% Le troisième la légende
% Le quatrième le label
Ici, je cite la figure \ref{fig:my_label} dans le texte.
\begin{figure}[h!]
    \centering
    \includegraphics[width=0.8\textwidth]{figures/vincent-brunie.jpg}
    \caption{Mettre une légende explicite à votre figure, ici M. Brunie Directeur de l'INSA Rennes.}
    \label{fig:my_label}
\end{figure}

\subsection{Insertion d'équation}
%------- Pour insérer et citer une équation --------------

\begin{equation} \label{eq: exemple}
\rho + \Delta = 42
\end{equation}

L'équation \ref{eq: exemple} est citée ici. 

\subsection{Insertion d'une référence bibliographique}
Les références (articles scientifiques, articles de journaux, blogs, pages web) doivent être mentionnées dans le texte par une balise \cite{maref} et fait le lien avec la citation incluse dans la bibliographie.

\bibliographystyle{apalike} % définit le style de bibliographie, ici APA
\bibliography{references} % ajoute la bibliographie basée sur le fichier references.bib



%----------- Résumé Français / Anglais -------------------
\newpage
\vspace*{\stretch{1}}
\begin{center}
	\begin{abstract}
        \lipsum[1]
        \newline
        \rule{\linewidth}{0.2 mm} \\[0.4 cm]
        \begin{center}\textbf{Summary :}\end{center} 
        \lipsum[1]
    \end{abstract}
\end{center}
\vspace*{\stretch{1}}


%------------ 4ème de couverture de l'INSA ---------------- 
\newpage
\includepdf [pages=-,noautoscale = true,scale=1] {documents/4e-couverture.pdf}


\end{document}

